% !TeX root = ActuarialFormSheet_MBFA-es.tex
% !TeX spellcheck = es_ES


\begin{f}[Esquemas de cesión/retrocesión]

\tikzstyle{Sources} = [rectangle, rounded corners, minimum width=5cm, minimum height=1cm,text centered, text width=5cm, draw=black, fill=BleuProfondIRA!40]
\tikzstyle{projections} = [ellipse, trapezium left angle=70, trapezium right angle=110, minimum height=1cm, text centered, draw=black, fill=FushiaIRA!30]
\tikzstyle{Calculs} = [rectangle, minimum width=5cm, minimum height=1cm, text centered, text width=5cm, draw=black, fill=OrangeProfondIRA!30]
\tikzstyle{modalite} = [ellipse, minimum width=3cm, minimum height=1cm, text centered, draw=black, fill=BleuProfondIRA!40]
\tikzstyle{arrow} = [thick,->,>=stealth]


\resizebox{\linewidth}{!}{			%
\begin{tikzpicture}[node distance=1.5cm]
\node (Assure) [Sources] {\begin{tabular}{c} ASEGURADO\\ \rowcolor{BleuProfondIRA!40}\tiny suscriptor\end{tabular}};
\node (Ca) [projections, below left = of Assure, xshift=2.5cm, yshift=.5cm] {Contrato de seguro};
\node (AG) [modalite, below right = of Assure, xshift=-2.5cm, yshift=.5cm] {Agente general / Corredor};
\node (Assureur) [Calculs, below of=Assure, node distance = 3cm] {\begin{tabular}{c}ASEGURADORA DIRECTA\\
\rowcolor{OrangeProfondIRA!30}  \tiny  cedente\end{tabular} };
\node (ConvR) [projections, below left = of Assureur, xshift=2.5cm, yshift=.5cm] {Convenio de reaseguro};
\node (Courtr) [modalite, below right = of Assureur, xshift=-2.5cm, yshift=.5cm] {Corredor de reaseguros};
\node (Reass) [Calculs, below of=Assureur, node distance = 3cm] {REASEGURADORA(S)};
\node (ConvR2) [projections, below left = of Reass, xshift=2.5cm, yshift=.5cm] {Convenio de reaseguro};
\node (Courtr2) [modalite, below right = of Reass, xshift=-2.5cm, yshift=.5cm] {Corredor de reaseguros};
\node (Retro) [Calculs, below of=Reass, node distance = 3cm] {RETROCESIONARIO(S)};
\draw [arrow] (Assure.south) -- (Assureur.north);
\draw [arrow] (Assureur.south) -- (Reass.north);
\draw [arrow] (Assureur.south) -- (Reass.north);
\draw [arrow] (Reass.south) -- (Retro.north);
%
\end{tikzpicture}}

\end{f}
\hrule


%https://www.reinsurancene.ws/top-50-reinsurance-groups/

\begin{f}[Términos clave del reaseguro]
	
	\textbf{Cedente:} cliente del reasegurador, es decir, la aseguradora directa, que transfiere (cede) riesgos al reasegurador a cambio del pago
	de una \textbf{prima de reaseguro}.
	
	\textbf{Cesión:} transferencia de riesgos por parte del asegurador directo al reasegurador.
	
	\textbf{Capacidad (\emph{Value Exposure})}: límite del importe del riesgo cubierto por un contrato de (re)seguro.
	
	\textbf{Reaseguro proporcional:} participación proporcional del reasegurador en las primas y en los siniestros del asegurador directo.
	
	\textbf{Reaseguro a cuota parte (\emph{Quota Share}):} tipo de reaseguro proporcional en el que el  reasegurador participa en un porcentaje determinado de todos los riesgos suscritos por un asegurador directo en un ramo concreto.
	
	\textbf{Reaseguro de excedente de retención (\emph{Surplus Share}):} tipo de reaseguro proporcional en el que el reasegurador cubre los riesgos que superan la retención
	del asegurador directo. Este ratio se calcula sobre la capacidad del riesgo suscrito (\(\approx\) siniestro máximo posible).
	
	\textbf{Comisión de reaseguro:} remuneración que el reasegurador concede al asegurador o a los corredores en compensación por los gastos de adquisición y gestión de los contratos de seguro.
	
	\textbf{Reaseguro no proporcional (o reaseguro de exceso de siniestro, \emph{Excess Reinsurance}):}
	asunción por parte del reasegurador de los siniestros que superen un importe determinado, a cambio del pago por parte del asegurador directo de una prima de reaseguro específica.
	
	\textbf{Retrocesión:}\index{R\'eassurance! Retrocesión}  parte de los riesgos que el reasegurador cede a otros reaseguradores.
	
	\textbf{Coaseguro:}\index{R\'eassurance! Coaseguro} participación de varios aseguradores directos en un mismo riesgo.
	
	En este caso se utiliza la expresión \textbf{pool de reaseguro}. 
	El reasegurador principal se denomina \textbf{apériteur}. 
	
	\textbf{Tratado de reaseguro:} contrato celebrado entre el asegurador directo y el reasegurador sobre una o varias carteras del asegurador.
	
	\textbf{Reaseguro facultativo:}
	Se diferencia del tratado de reaseguro en que la suscripción se realiza riesgo por riesgo (o póliza por póliza) (caso por caso, un riesgo cada vez). 
	
\end{f}
\hrule




\begin{f}
	[El papel económico del reaseguro]

El seguro y el reaseguro comparten el mismo objetivo: la mutualización de los riesgos.
El reaseguro interviene, en particular, en los siguientes riesgos:
\begin{itemize}
	\item independientes, pero costosos individualmente (aviones, buques, instalaciones industriales, etc.),
	\item de pequeña cuantía (roturas, automóviles, etc.), pero que se correlacionan en caso de sucesos de gran envergadura, lo que genera acumulaciones onerosas,
	\item agregados dentro de una cartera de pólizas de seguro, 
	\item poco conocidos o nuevos.
\end{itemize}


El reaseguro permite aumentar la capacidad de suscripción de negocios, garantizar la estabilidad financiera de la aseguradora, sobre todo en caso de catástrofes, reducir sus necesidades de capital y beneficiarse de la experiencia del reasegurador.

\end{f}	
\hrule



\begin{f}
	[Tipos de acuerdos de reaseguro]

{\color{white}.}	

\begin{center}
		\resizebox{0.90\linewidth}{!}{\input{Graph/TypeAgreements-es}}
\end{center}
\end{f}
\hrule


\begin{f}[Tipos de reaseguro a través de un ejemplo]
	
	Nuestra aseguradora reasegura \(N=30\) pólizas de seguro, con un total de primas de 10M\EUR{} (\(P=\sum_{i=1\ldots N}P_i\)).
	La capacidad total es de 180M\EUR{} (\(\sum_{i=1\ldots30}K_i\)).
	\(S_r\) será la parte total del siniestro asumida por la aseguradora y \(P_r\) la prima total de reaseguro.
	A continuación se presentan las \(n=8\) pólizas con siniestros (\(1\geq i \geq n\)), siendo los siniestros de las demás pólizas nulos (\(S_i=0, \forall i>n\)):
	
	\begin{center}\footnotesize
		\renewcommand{\arraystretch}{1.25}
		\begin{tabular}{|l|rrrrrrrr|}
			\hline
			\rowcolor{BleuProfondIRA!40}         Número de siniestro    & 1 & 2 & 3 & 4 & 5 & 6 & 7 & 8 \\ \hline 
			Prima  (k\EUR{})        & 500 & 200 & 100 & 100 & 50 & 200 & 500 & 200 \\ 
			Capacidad  (M\EUR{})       & 8 & 5 & 3 & 2 & 3 & 5 & 8 & 8 \\ 
			Siniestros  (M\EUR{})       & 1 & 1 & 1 & 2 & 3 & 3 & 5 & 8 \\ \hline
		\end{tabular}
		\renewcommand{\arraystretch}{1}
	\end{center}
	
	El \(S/P\) es del 240\%.
	
	
	\begin{center}
		\begin{tikzpicture}
			\begin{axis}[axis x line=bottom, axis y line = left,ymin=0,ymax=24,xmin=0, xmax=10, height=7cm,width=5cm, xticklabel style ={font=\footnotesize,align=center},  yticklabel style ={font=\footnotesize}, 
				legend style={at={(0.5,-0.15)},anchor=north,legend columns=-1,font=\tiny}]
				\addplot [ybar,fill=BleuProfondIRA!40,mark=none,draw=none] table [x index=0, y index=1] {..\\_Common\\ReinsuranceClaim.dat} ;
				\legend{Siniestros tramitados}
			\end{axis}	
		\end{tikzpicture}
		\begin{tikzpicture}
			\begin{axis}[axis x line=center,axis y line = left,ymin=0,ymax=24,xmin=0, xmax=10,height=7cm,width=5cm, xticklabel style ={font=\footnotesize,align=center},  yticklabel style ={font=\footnotesize}, 
				legend style={at={(0.5,-0.15)},
					anchor=north,legend columns=-1,font=\tiny}]
				\addplot [ybar,fill=BleuProfondIRA!40,mark=none,draw=none] table [x index=0, y index=2] {..\\_Common\\ReinsuranceClaim.dat} ;
				\legend{Siniestros acumulados}
			\end{axis}	
		\end{tikzpicture}
		%\includegraphics{../../Reinsurance_M2/Graph/ReinsuranceClaim.pdf}
		%\includegraphics{../../Reinsurance_M2/Graph/ReinsuranceClaimCum.pdf}
	\end{center}
\medskip	


\textbf{Cuota:}
\[
S_r=\alpha \sum_{i=1\ldots n}S_i\quad \ \ P_r=\alpha\sum_{i=1\ldots N}P_i 
\]
donde \(\alpha\ \in [0,1]\) (25 \% en la figura) es la parte cedida en concepto de cuota.
	
	%\setlength{\tabcolsep}{1cm}

%\begin{tabular}{|l|rrrrrrrr|}
%	\hline
%	\rowcolor{BleuProfondIRA!40}         Num de sinistre		& 1 & 2 & 3 & 4 & 5 & 6 & 7 & 8 \\ \hline \hline
%	Sinistres  (M\EUR{})   	& 1 & 1 & 1 & 2 & 3 & 3 & 5 & 8 \\ 
%	Cédante  (M\EUR{})   	& 0,75 & 0,75 & 0,75 &
%	1,50 & 2,25 & 2,25 & 
%	3,75 & 6,00\\
%	Réassurance  (M\EUR{})  & 0,25 & 0,25 & 0,25 &
%	0,50 & 0,75 & 0,75 & 
%	1,25 & 2,00\\
%	Prime portef conservée  & \multicolumn{2}{r}{\cellcolor{mbfaulmbleu!10}7,5 M\EUR{}} & & & & & &\\
%	Sinistre conservé  	& \multicolumn{2}{r}{18} & & & & & & \\
%	S/P conservé  		& \multicolumn{2}{r}{\cellcolor{mbfaulmbleu!10}240\%} & & & & & & \\
%	\hline
%\end{tabular}
	
\begin{center}
		\begin{tikzpicture}
		\begin{axis}[axis x line=bottom, axis y line = left,ymin=0,ymax=24,xmin=0, xmax=10, height=7cm,width=5cm, xticklabel style ={font=\footnotesize,align=center},  yticklabel style ={font=\footnotesize}, 
			legend style={at={(0.5,-0.15)},anchor=north,legend columns=1,font=\tiny}]
			\addplot [ybar ,fill=OrangeMoyenIRA!40,mark=none,draw=OrangeMoyenIRA!40] table [x index=0, y index=1] {..\\_Common\\ReinsuranceClaim.dat} ;
			\addplot [ybar ,fill=BleuProfondIRA!40,mark=none,draw=BleuProfondIRA!40] table [x index=0, y index=6] {..\\_Common\\ReinsuranceClaim.dat} ;
			\legend{Siniestros cedidos, siniestros retenidos QP}
		\end{axis}	
	\end{tikzpicture}
	%	
	\begin{tikzpicture}
		\begin{axis}[axis x line=center,axis y line = left,ymin=0,ymax=24,xmin=0, xmax=10,height=7cm,width=5cm, xticklabel style ={font=\footnotesize,align=center},  yticklabel style ={font=\footnotesize}, 
			legend style={at={(0.5,-0.15)},
				anchor=north,legend columns=1,font=\tiny}]
			\addplot [ybar ,fill=OrangeMoyenIRA!40,mark=none,draw=OrangeMoyenIRA!40] table [x index=0, y index=2] {..\\_Common\\ReinsuranceClaim.dat} ;
			\addplot [ybar ,fill=BleuProfondIRA!40,mark=none,draw=BleuProfondIRA!40] table [x index=0, y index=7] {..\\_Common\\ReinsuranceClaim.dat} ;
			\legend{Siniestros cedidos, Siniestros retenidos QP}
		\end{axis}	
	\end{tikzpicture}
	%\includegraphics{../../Reinsurance_M2/Graph/ReinsuranceClaimQuotePart.pdf}
	%\includegraphics{../../Reinsurance_M2/Graph/ReinsuranceClaimCumQuotePart.pdf}

\end{center}
\medskip

\textbf{Excédent de plein}, le plein est noté \(\boldsymbol{K}\) (2M\EUR{} dans l'exemple), \(\alpha_i\) représente le taux de cession de la police \(i\).
	\[
	S_r= \sum_{i=1\ldots n}\underbrace{\left(\frac{\left(K_i-\boldsymbol{K} \right)_+ }{K_i} \right)}_{\alpha_i} S_i\quad \ \ P_r=\sum_{i=1\ldots N}\left(\frac{\left(K_i-\boldsymbol{K} \right)_+ }{K_i} \right)P_i 
	\]
	%\renewcommand{\arraystretch}{1.25}
	%\rowcolors{1}{sectionColor}{white}
	%\setlength{\tabcolsep}{1cm}
%\begin{tabular}{|l|rrrrrrrr|}
%\hline
%\rowcolor{BleuProfondIRA!40}       Num de sinistre		& 1 & 2 & 3 & 4 & 5 & 6 & 7 & 8 \\ \hline \hline
%Sinistres  (M\EUR{})   	& 1 & 1 & 1 & 2 & 3 & 3 & 5 & 8 \\ 
%Capacité  (M\EUR{})   	& 8 & 5 & 3 & 2 & 3 & 5 & 8 & 8 \\ 
%Capacité cédante	& \multicolumn{2}{r}{60M\EUR{}} & & & & & & \\
%Prime cédante		& \multicolumn{2}{r}{\cellcolor{mbfaulmbleu!40}3M\EUR{}} & & & & & & \\
%Part cédée  (\%)   & 75 & 60 & 33 &
%0 & 33 & 60 & 
%75 & 75\\
%Cédante  (M\EUR{})  & 0,25 & 0,4 & 0,67 &
%2,00 & 2,00 & 1,20 & 
%1,25 & 2,00\\
%Réassurance  (M\EUR{})  & 0,75 & 0,60 & 0,33 &
%0,00 & 1,00 & 1,80 & 
%3,75 & 6,00\\
%Sinistre conservé  	& \multicolumn{2}{r}{\cellcolor{mbfaulmbleu!40}9,77} & & & & & & \\
%Sinistre cédé  		& \multicolumn{2}{r}{14,23} & & & & & & \\
%S/P conservée  		& \multicolumn{2}{r}{\cellcolor{mbfaulmbleu!40}325\%} & & & & & & \\
%\hline
%\end{tabular}


\begin{center}
		%\includegraphics{../../Reinsurance_M2/Graph/ReinsuranceClaimEPlein.pdf}
	%\includegraphics{../../Reinsurance_M2/Graph/ReinsuranceClaimCumEPlein.pdf}
	\begin{tikzpicture}
		%format de la date yyyy-mm-dd , pas de nom de colonne vide,
		\begin{axis}[axis x line=bottom, axis y line = left,ymin=0,ymax=24,xmin=0, xmax=10, height=7cm,width=5cm, xticklabel style ={font=\footnotesize,align=center},  yticklabel style ={font=\footnotesize}, 
			legend style={at={(0.5,-0.15)},anchor=north,legend columns=2,font=\tiny}]
			\addplot [ybar ,mark=none,draw=OrangeMoyenIRA!40,fill=OrangeMoyenIRA!40] table [x index=0, y index=1] {..\\_Common\\ReinsuranceClaim.dat} ;
			\addplot [ybar ,mark=none,draw=GrisLogoIRA,mark=none] table [x index=0, y index=11] {..\\_Common\\ReinsuranceClaim.dat} ;
			\addplot [draw=GrisLogoIRA] table [x index=0, y index=4] {..\\_Common\\ReinsuranceClaim.dat} ;
			\addplot [ybar ,fill=BleuProfondIRA!40,mark=none,draw=BleuProfondIRA!40] table [x index=0, y index=12] {..\\_Common\\ReinsuranceClaim.dat} ;
			\legend{Sin cedidos, Capacidad, Lleno, Sin guardados}
		\end{axis}	
	\end{tikzpicture}
	%	
	\begin{tikzpicture}
		%format de la date yyyy-mm-dd , pas de nom de colonne vide,
		\begin{axis}[axis x line=center,axis y line = left,ymin=0,ymax=24,xmin=0, xmax=10,height=7cm,width=5cm, xticklabel style ={font=\footnotesize,align=center},  yticklabel style ={font=\footnotesize}, 
			legend style={at={(0.5,-0.15)},
				anchor=north,legend columns=1,font=\tiny}]
			\addplot [ybar ,fill=OrangeMoyenIRA!40,mark=none,draw=OrangeMoyenIRA!40] table [x index=0, y index=2] {..\\_Common\\ReinsuranceClaim.dat} ;
			\addplot [ybar ,fill=BleuProfondIRA!40,mark=none,draw=BleuProfondIRA!40] table [x index=0, y index=13] {..\\_Common\\ReinsuranceClaim.dat} ;
			%		\addplot [ybar interval,fill=BleuProfondIRA,mark=none,draw=BleuProfondIRA] table [x index=0, y index=7] {..\\_Common\\ReinsuranceClaim.dat} ;
			\legend{Siniestros cedidos, Siniestros retenidos}
		\end{axis}	
	\end{tikzpicture}
\end{center}
\medskip

\textbf{Franquicia por siniestro}
	%\renewcommand{\arraystretch}{1.25}
	%\rowcolors{1}{sectionColor}{white}
	%\setlength{\tabcolsep}{1cm}
	
La aseguradora fija la prioridad \(a\) y el alcance \(b\) (2M\EUR{} y 4M\EUR{}, respectivamente, en la figura).
\[
S_r= \sum_{i=1\ldots n} \min\left( \left( S_i-a\right)^+,b\right)  
\]
La prima la fija la reaseguradora, en función de su estimación de \(\mathbb{E}[S_r]\).
%\begin{tabular}{|l|rrrrrrrr|}
%	\hline
%	\rowcolor{BleuProfondIRA!40}        Num de sinistre		& 1 & 2 & 3 & 4 & 5 & 6 & 7 & 8 \\ \hline \hline
%	Sinistres  (M\EUR{})   	& 1 & 1 & 1 & 2 & 3 & 3 & 5 & 8 \\ 
%	Cédante  (M\EUR{})   	& 1 & 1 & 1 &
%	2 & 2 & 2 & 
%	2 & 4\\
%	Réassurance  (M\EUR{})  & 0 & 0 & 0 &
%	0 & 1 & 1 & 
%	3 & 4\\
%	\hline
%\end{tabular}

\begin{center}
		%\includegraphics{../../Reinsurance_M2/Graph/ReinsuranceClaimXSsin.pdf}
	%\includegraphics{../../Reinsurance_M2/Graph/ReinsuranceClaimCumXSsin.pdf}
	\begin{tikzpicture}
		%format de la date yyyy-mm-dd , pas de nom de colonne vide,
		\begin{axis}[axis x line=bottom, axis y line = left,ymin=0,ymax=24,xmin=0, xmax=10, height=7cm,width=5cm, xticklabel style ={font=\footnotesize,align=center},  yticklabel style ={font=\footnotesize}, 
			legend style={at={(0.5,-0.15)},anchor=north,legend columns=2,font=\tiny}]
			\addplot [ybar ,mark=none,draw=BleuProfondIRA!40,fill=BleuProfondIRA!40] table [x index=0, y index=1] {..\\_Common\\ReinsuranceClaim.dat} ;
			\addplot [ybar ,mark=none,draw=OrangeMoyenIRA!40,fill=OrangeMoyenIRA!40] table [x index=0, y index=10] {..\\_Common\\ReinsuranceClaim.dat} ;
			\addplot [ybar ,mark=none,draw=BleuProfondIRA!40,fill=BleuProfondIRA!40] table [x index=0, y index=8] {..\\_Common\\ReinsuranceClaim.dat} ;
			\addplot [draw=GrisLogoIRA] table [x index=0, y index=4] {..\\_Common\\ReinsuranceClaim.dat} ;
			\addplot [draw=GrisLogoIRA] table [x index=0, y index=5] {..\\_Common\\ReinsuranceClaim.dat} ;
			\legend{Sin conservados, sin cedidos, prioridad 2 M€, alcance 4 M€ }
		\end{axis}	
	\end{tikzpicture}
	%	
	\begin{tikzpicture}
		%format de la date yyyy-mm-dd , pas de nom de colonne vide,
		\begin{axis}[axis x line=center,axis y line = left,ymin=0,ymax=24,xmin=0, xmax=10,height=7cm,width=5cm, xticklabel style ={font=\footnotesize,align=center},  yticklabel style ={font=\footnotesize}, 
			legend style={at={(0.5,-0.15)},
				anchor=north,legend columns=1,font=\tiny}]
			\addplot [ybar ,fill=OrangeMoyenIRA!40,mark=none,draw=OrangeMoyenIRA!40] table [x index=0, y index=2] {..\\_Common\\ReinsuranceClaim.dat} ;
			\addplot [ybar ,fill=BleuProfondIRA!40,mark=none,draw=BleuProfondIRA!40] table [x index=0, y index=9] {..\\_Common\\ReinsuranceClaim.dat} ;
			%		\addplot [ybar ,fill=BleuProfondIRA,mark=none,draw=BleuProfondIRA] table [x index=0, y index=7] {..\\_Common\\ReinsuranceClaim.dat} ;
			\legend{Siniestros cedidos, Siniestros retenidos}
		\end{axis}	
	\end{tikzpicture}
	
	\emph {WXL-R = Working XL per Risk}
\end{center}
	\medskip


\textbf{Franquicia por siniestro}
	%\renewcommand{\arraystretch}{1.25}
	%\rowcolors{1}{sectionColor}{white}
	%\setlength{\tabcolsep}{1cm}
	
	\[
	S_r= \sum_{\begin{array}{c}
			Cat_j,\\ i=1\ldots N
	\end{array}} \min\left( \mathds{1}_{i\in Cat_j}\times \left( S_i-a\right)^+,b\right)  
	\]
	
	En la ilustración, los siniestros se refieren a un único suceso, con una prioridad de 5M\EUR{} y un alcance de 10M\EUR{}.
	
%\begin{tabular}{|l|rrrrrrrr|}
%	\hline
%	\rowcolor{BleuProfondIRA!40}        Num de sinistre		& 1 & 2 & 3 & 4 & 5 & 6 & 7 & 8 \\ \hline \hline
%	Sinistres  (M\EUR{})   	& 1 & 1 & 1 & 2 & 3 & 3 & 5 & 8 \\ 
%	Sinistres en Cumulé   	& 1 & 2 & 3 & 5 & 8 & 11 & 16 & 24 \\ 
%	Cédante en Cumulé   	& 1 & 2 & 3 &
%	5 & 5 & 5 & 
%	6 & 14\\
%	Réassurance en Cumulé  & 0 & 0 & 0 &
%	0 & 3 & 6 & 
%	10 & 10\\
%	\hline
%\end{tabular}

	
	%{../../Reinsurance_M2/Graph/ReinsuranceClaimXScat.pdf}
	%\includegraphics{../../Reinsurance_M2/Graph/ReinsuranceClaimCumXScat.pdf}
\begin{center}
		\begin{tikzpicture}
		\begin{axis}[axis x line=center,axis y line = left,ymin=0,ymax=24,xmin=0, xmax=10,height=7cm,width=5cm, xticklabel style ={font=\footnotesize,align=center},  yticklabel style ={font=\footnotesize}, 
			legend style={at={(0.5,-0.15)},
				anchor=north,legend columns=1,font=\tiny}]
			\addplot [ybar ,fill=OrangeMoyenIRA!40,mark=none,draw=OrangeMoyenIRA!40] table [x index=0, y index=1] {..\\_Common\\ReinsuranceClaim.dat} ;
			\addplot [ybar ,fill=BleuProfondIRA!40,mark=none,draw=BleuProfondIRA!40] table [x index=0, y index=18] {..\\_Common\\ReinsuranceClaim.dat} ;
			\legend{Siniestros cedidos, Siniestros retenidos};
		\end{axis}	
	\end{tikzpicture}
	%	
	\begin{tikzpicture}
		\begin{axis}[axis x line=bottom, axis y line = left,ymin=0,ymax=24,xmin=0, xmax=10, height=7cm,width=5cm, xticklabel style ={font=\footnotesize,align=center},  yticklabel style ={font=\footnotesize}, 
			legend style={at={(0.5,-0.15)},anchor=north,legend columns=2,font=\tiny}]
			\addplot [ybar ,mark=none,draw=BleuProfondIRA!40,fill=BleuProfondIRA!40] table [x index=0, y index=2] {..\\_Common\\ReinsuranceClaim.dat} ;
			\addplot [ybar ,mark=none,draw=OrangeMoyenIRA!40,fill=OrangeMoyenIRA!40] table [x index=0, y index=16] {..\\_Common\\ReinsuranceClaim.dat} ;
			\addplot [ybar ,mark=none,draw=BleuProfondIRA!40,fill=BleuProfondIRA!40] table [x index=0, y index=17] {..\\_Common\\ReinsuranceClaim.dat} ;
			\addplot [draw=GrisLogoIRA] table [x index=0, y index=14] {..\\_Common\\ReinsuranceClaim.dat} ;
			\addplot [draw=GrisLogoIRA] table [x index=0, y index=15] {..\\_Common\\ReinsuranceClaim.dat} ;
			\legend{Sin conservados, sin cedidos,, prioridad 2 M€, alcance 4 M€ };
		\end{axis}	
	\end{tikzpicture}
	
	\emph{ Cat-XL = Catastrophe XL }
\end{center}
\medskip

	
	
	%\renewcommand{\arraystretch}{1.25}
	%\rowcolors{1}{sectionColor}{white}
	%\setlength{\tabcolsep}{1cm}
	
%	L'assureur fixe la priorité à 5M\EUR{}  et la portée à 10M\EUR{}, soit un plafond de 15M\EUR{}.
%	Dans notre exemple supposons que les 8 sinistres font référence à \textbf{deux événements, une première tempête en octobre, une deuxième en décembre}.
	
%\begin{tabular}{|l|rrrrrrrr|}
%	\hline
%	\rowcolor{BleuProfondIRA!40}         Num de sinistre		& 1 & 2 & 3 & 4 & 5 & 6 & 7 & 8 \\ \hline \hline
%	Événement		&Oct & Oct & Déc & Oct & Déc & Oct & Déc & Déc \\ \hline \hline
%	Sinistres  (M\EUR{})   	& 1 & 1 & 1 & 2 & 3 & 3 & 5 & 8 \\ 
%	\hline
%\end{tabular}
%\begin{tabular}{|l|rr|r|}
%	\hline
%	\rowcolor{BleuProfondIRA!40}        Événement		&Oct &  Déc  & Total\\ \hline \hline
%	Som de sinistre		& 7 & 17& 24\\ \hline \hline
%	Cédante  	& 5 & 7& 12 \\
%	Réassurance en Cumulé  & 2 & 10& 12 \\
%	\hline
%\end{tabular}



\textbf{Excédente de pérdidas anuales}
	
	
Este reaseguro (\emph{Stop Loss}) se activa cuando el total de pérdidas anuales se deteriora. 
Se expresa en función del ratio \(S/P\), con una prioridad y un alcance del \(XL\) expresados en \%.
	
\[
S_r= \min\left( \left( \sum_{i=1\ldots N} S_i - a  P\right)^+,bP\right)  
\]

\end{f}
\hrule
	
\begin{f}[Las principales cláusulas en materia de reaseguro]
	
	\textbf{La franquicia} \(a^{ag}\) y el \textbf{límite agregado} \(b^{ag}\) se aplican tras el cálculo de \(S_r\). 
	
\[
S_r^{ag} = \min\left( \left(S_r - a^{ag}  \right)^+,b^{ag}\right)  
\]
\medskip

El objetivo de la \textbf{cláusula de indexación} es mantener las \underline{condiciones del contrato} a lo largo de varios ejercicios consecutivos.  Los límites del contrato se ajustan a un índice económico (salario, divisa, índice de precios, etc.).
\medskip

Con la \textbf{cláusula de estabilización}, cuando la siniestralidad se ve afectada por una \underline{liquidación prolongada}, o incluso muy prolongada (al menos \(\geq 1\) año), los límites del contrato se actualizan en el cálculo de \(S_r\) para que se respeten, en términos generales, las participaciones respectivas del reasegurador y de la cedente previstas inicialmente.
\medskip


Con la \textbf{cláusula de reparto de intereses}, si en una transacción o en una sentencia judicial se ha establecido una distinción \underline{entre
	la indemnización y los intereses}, los intereses devengados entre la fecha del siniestro y la del pago efectivo de la indemnización se repartirán entre la cedente y el reasegurador proporcionalmente a
su carga respectiva resultante de la aplicación del tratado, sin incluir los intereses.
\medskip


La \textbf{cláusula de reconstitución de la garantía}
se refiere únicamente a los \emph{tratados de exceso de siniestro por riesgo o por evento} que podrían activarse varias veces al año.
El reasegurador \underline{limita su prestación a \(N \) veces el alcance} del \(XS\), a cambio del pago de una prima complementaria. 
La reconstitución puede realizarse a prorrata temporis (tiempo que queda por transcurrir hasta la fecha
de vencimiento del tratado) o a prorrata de los capitales absorbidos, o ambas (doble prorrata). 

La \textbf{cláusula de superposición}
(\emph{Interlocking Clause})
se utiliza en los tratados en \(XS\) por evento,
que funcionan por ejercicio de suscripción y no \underline{por ejercicio de ocurrencia}.
La cláusula de
superposición tendrá como efecto recalcular los límites del contrato, ya que un mismo evento puede activar los contratos de suscripción \(n\) y \(N-1\).

\end{f}
\hrule

\begin{f}[El reaseguro público]
	La \href{http://www.ccr.fr}{Caisse Centrale de Réassurance (CCR)} ofrece, con la garantía del Estado, coberturas ilimitadas para ramos específicos del mercado francés.
	\begin{itemize}
		\item   los riesgos excepcionales relacionados con el transporte,
		\item   la responsabilidad civil de los operadores de buques e instalaciones nucleares,
		\item   los riesgos de catástrofes naturales,
		\item   los riesgos de atentados y actos de terrorismo,
		\item   el Complemento de Seguro de Crédito Público (CAP).
		\end {itemize}
		Asimismo, gestiona en nombre del Estado determinados fondos públicos, en particular el régimen Cat Nat.
		
		Asimismo, la \href{http://www.gareat.com/}{GAREAT} es una Agrupación de Interés Económico (GIE) sin ánimo de lucro, mandatada por sus miembros, que gestiona el reaseguro de los riesgos de atentados y actos de terrorismo con el apoyo del Estado a través de la CCR.
\end{f}
\hrule

\begin{f}[Titulización / CatBonds]
	\label {CatBond}
	¿Por qué? La capacidad financiera conjunta de todas las aseguradoras y reaseguradoras no  alcanza para cubrir los daños de un terremoto de gran magnitud en Estados Unidos (\(\geq\) 200 mil millones de euros{}). 
	Esta suma representa  menos del 1\% de la capitalización de los mercados financieros estadounidenses. 
	
	La \textbf{titulización} transforma un riesgo asegurador en un título negociable, a menudo en títulos de deuda denominados \textbf{Cat-Bonds}. 
	Consiste en un intercambio de capital a cambio del pago periódico de cupones, en el que el pago de los cupones y/o el reembolso del capital están condicionados a la ocurrencia de un evento desencadenante definido \emph{a priori}.
	Los tipos de interés de estos bonos se incrementan en función del riesgo, no de impago o de contraparte, sino de la ocurrencia del evento (inferior al 1\%). La estructura dedicada a esta transformación se denomina \emph{Special Purpose Vehicle} (SPV).
	
	
	El evento desencadenante puede estar directamente vinculado a los resultados de la entidad cedente (indemnización), depender de un índice de siniestralidad, de un parámetro medible (suma de los excedentes de lluvia, escala de Richter, tasa de mortalidad) o de un modelo (RMS y Equecat Storm modelling). 
	

	\renewcommand{\arraystretch}{1.25}
\begin{center}\small
\begin{tabular}{|m{20mm}|*{4}{>{\centering\arraybackslash}m{14mm}|}}
\hline \rowcolor{BleuProfondIRA!40}
\textbf{Criterio} & \textbf{Indemni\-zación} & \textbf{Índice} & \textbf{Pará\-metro} & \textbf{Modelo} \\
\hline
Transparencia & \(\ominus\) & \(\oplus\) & \(\oplus\) & \(\oplus\) \\
\hline
Riesgo básico & \(\oplus\) & \(\ominus\) & \(\ominus\) & \(\oplus\) \\
\hline
Riesgo moral & \(\ominus\) & \(\oplus\) & \(\oplus\) & \(\oplus\) \\
\hline
Universalidad de los riesgos & \(\oplus\) & \(\oplus\) & \(\ominus\) & \(\oplus\) \\
\hline
Plazo de activación & \(\ominus\) & \(\ominus\) & \(\oplus\) & \(\oplus\) \\
\hline
\end{tabular}
\end{center}	\renewcommand{\arraystretch}{1}
\end{f}






